\chapter{About This Software}

\section{What problem does it solve?}

A large denim factory makes tens of thousands of garments every day. Those
garments are constantly moving: from the cutting room to the sewing lines, into
the washing plant, on to finishing, through quality checking, and finally into
cartons for the customer.

The old way of keeping track was paper. Somebody counted a bundle, wrote the
number on a slip, and passed it along. When the numbers did not match at the end
of the day, nobody could say where the missing pieces went.

This software replaces that guesswork. A small electronic tag --- called an
\textbf{RFID tag} --- is attached to every single garment. Special readers can
sense hundreds of these tags at once, without anyone having to look at them or
scan them one by one.

\begin{plainwords}
\textbf{RFID} stands for ``Radio Frequency Identification''. Think of it as a
tiny electronic name badge. It has no battery. When a reader sends out a radio
signal, the tag wakes up and calls back its own unique number. A reader can hear
hundreds of tags in a couple of seconds, even inside a closed trolley.
\end{plainwords}

\section{What the system keeps track of}

\begin{bullets}
  \item \textbf{Fabric rolls} arriving from suppliers, and how much of each roll has been used.
  \item \textbf{Bundles} of cut pieces going from the cutting room to the sewing lines.
  \item \textbf{Every finished garment}, individually, from stitching to shipping.
  \item \textbf{Every handover} between departments, with a document and a count check.
  \item \textbf{Every quality inspection}, including exactly where on the garment a fault was found.
  \item \textbf{Every person} who touched the work, and when.
\end{bullets}

\section{The journey of one pair of jeans}

Figure~\ref{fig:flow} shows the whole route. Each box is a department. Each
arrow is a handover, and at every handover the system does the same two things:
the department sending the goods scans what is leaving, and the department
receiving them scans what arrived. The computer then compares the two.

\begin{figure}[H]
\centering
\begin{tikzpicture}[
  node distance=6mm and 9mm,
  box/.style={rectangle,rounded corners=3pt,draw=none,text=white,
              minimum height=11mm,minimum width=27mm,align=center,font=\small\bfseries},
  arrow/.style={-{Stealth[length=3mm]},thick,draw=brandgrey},
  lbl/.style={font=\scriptsize\itshape,text=brandgrey}]

\node[box,fill=secFabric] (fab) {Fabric\\Warehouse};
\node[box,fill=secCut,right=of fab] (cut) {Cutting};
\node[box,fill=secStitch,right=of cut] (sti) {Stitching};
\node[box,fill=secSort,right=of sti] (sor) {Sorting};

\node[box,fill=secWash,below=14mm of sor] (was) {Washing \&\\Treatment};
\node[box,fill=secFinish,left=of was] (fin) {Finishing};
\node[box,fill=secQC,left=of fin] (qc) {Quality\\Control};
\node[box,fill=secRetro,left=of qc] (ret) {Retrofitting};

\node[box,fill=secDisp,below=14mm of ret] (dis) {Dispatch \&\\Packing};
\node[box,fill=brandgrey,right=of dis] (shp) {Shipped};

\draw[arrow] (fab) -- node[lbl,above]{rolls} (cut);
\draw[arrow] (cut) -- node[lbl,above,align=center]{hand count} (sti);
\draw[arrow] (sti) -- node[lbl,above,align=center]{tag added} (sor);
\draw[arrow] (sor) -- (was);
\draw[arrow] (was) -- (fin);
\draw[arrow] (fin) -- (qc);
\draw[arrow] (qc) -- node[lbl,above]{failed} (ret);
\draw[arrow] (ret) -- (dis);
\draw[arrow] (dis) -- node[lbl,above,align=center]{new tag} (shp);

\draw[arrow,dashed] (ret.north) to[out=60,in=200] node[lbl,above,sloped]{corrected, back to QC} (qc.south west);
\draw[arrow,dashed] (qc.south) to[out=280,in=20] node[lbl,below,sloped]{passed} (dis.north east);

\end{tikzpicture}
\caption{The route every garment takes. Solid arrows are the normal path;
dashed arrows show what happens after a quality check.}
\label{fig:flow}
\end{figure}

\section{The one idea you must understand: the two-sided count}

This is the heart of the system, and it is simple.

\begin{figure}[H]
\centering
\begin{tikzpicture}[
  node distance=10mm,
  dept/.style={rectangle,rounded corners=3pt,fill=brandsoft,draw=brandblue,
               minimum height=16mm,minimum width=38mm,align=center,font=\small\bfseries},
  doc/.style={rectangle,rounded corners=2pt,fill=white,draw=brandgrey,
              minimum height=13mm,minimum width=34mm,align=center,font=\scriptsize},
  arrow/.style={-{Stealth[length=3mm]},thick,draw=brandblue}]

\node[dept] (a) {Sending\\department};
\node[doc,right=22mm of a] (d) {Transfer note\\\textbf{240 pieces}};
\node[dept,right=22mm of d] (b) {Receiving\\department};

\draw[arrow] (a) -- node[font=\scriptsize,above,align=center]{scans what\\is leaving} (d);
\draw[arrow] (d) -- node[font=\scriptsize,above,align=center]{scans what\\arrived} (b);

\node[below=9mm of d,font=\small,align=center,text=brandink]
  {The computer compares the two numbers};
\node[below=17mm of d,align=center,font=\small\bfseries]
  {\textcolor{okgreen}{Same $\Rightarrow$ MATCHED}\qquad
   \textcolor{dangerred}{Different $\Rightarrow$ VARIANCE}};
\end{tikzpicture}
\caption{Every handover between departments works this way.}
\label{fig:tally}
\end{figure}

If the two counts agree, the handover is marked \textbf{MATCHED} and everyone
moves on. If they do not agree, the system marks it a \textbf{VARIANCE} and
shows a list of exactly which garments were not found --- by serial number, style
and size. Nobody has to hunt through paperwork.

\begin{tipbox}
A variance is not an accusation. Tags are sometimes missed because a garment was
at the bottom of a trolley. Simply scan the batch again; anything found the
second time is added, and the variance clears by itself.
\end{tipbox}

\section{Who uses which portal}

The same software shows a different set of screens depending on who signs in.
This manual has one chapter for each group.

\begin{table}[H]
\centering
\small
\begin{tabular}{@{}p{34mm}p{45mm}p{60mm}@{}}
\toprule
\textbf{Portal} & \textbf{Who signs in} & \textbf{What they mainly do} \\
\midrule
Operator &
  Machine and station staff in the store, cutting, sewing, sorting, wash,
  finishing, retrofit and dispatch &
  Scan garments in and out, attach tags, count bundles, correct faults \\
Quality Inspector &
  QC staff at the inspection benches &
  Check garments, mark faults on the picture of the design, pass or fail \\
Supervisor &
  The person in charge of a department &
  Everything an operator can do, plus settle count differences and cancel
  mistaken handovers \\
Administrator &
  Plant management and the IT administrator &
  See all the numbers, build reports, manage staff accounts, register readers,
  maintain the product list \\
\bottomrule
\end{tabular}
\caption{The four portals covered by this manual.}
\end{table}
